\documentclass{report}
\usepackage{luatexja} % LuaTeXで日本語を使うためのパッケージ
\usepackage{luatexja-fontspec} % LuaTeX用の日本語フォント設定

% --- 数学関連 ---
\usepackage{amsmath, amssymb, amsfonts, mathtools, bm, amsthm} % 基本的な数学パッケージ
\usepackage{type1cm, upgreek} % 数式フォントとギリシャ文字
\usepackage{physics, mhchem} % 物理や化学の記号や式の表記を簡単にする

% --- 表関連 ---
\usepackage{multirow, longtable, tabularx, array, colortbl, dcolumn, diagbox} % 表のレイアウトを柔軟にする
\usepackage{tablefootnote, truthtable} % 表中に注釈を追加、真理値表
\usepackage{tabularray} % 高度な表組みレイアウト

% --- グラフィック関連 ---
\usepackage{tikz, graphicx} % 図の描画と画像の挿入
\usepackage{background} % ウォーターマークの設定
\usepackage{caption, subcaption} % 図や表のキャプション設定
\usepackage{float, here} % 図や表の位置指定

% --- レイアウトとページ設定 ---
\usepackage{fancyhdr} % ページヘッダー、フッター、余白の設定
\usepackage[top = 20truemm, bottom = 20truemm, left = 20truemm, right = 20truemm]{geometry}
\usepackage{fancybox, ascmac} % ボックスのデザイン

% --- 色とスタイル ---
\usepackage{xcolor, color, colortbl, tcolorbox} % 色とカラーボックス
\usepackage{listings, jvlisting} % コードの色付けとフォーマット

% --- 参考文献関連 ---
\usepackage{biblatex, usebib} % 参考文献の管理と挿入
\usepackage{url, hyperref} % URLとリンクの設定

% --- その他の便利なパッケージ ---
\usepackage{footmisc} % 脚注のカスタマイズ
\usepackage{multicol} % 複数段組
\usepackage{comment} % コメントアウトの拡張
\usepackage{siunitx} % 単位の表記
\usepackage{docmute}
% \usepackage{appendix}
% --- tcolorboxとtikzの設定 ---
\tcbuselibrary{theorems, breakable} % 定理のボックスと改ページ設定
\usetikzlibrary{decorations.markings, arrows.meta, calc} % tikzの装飾や矢印の設定

% --- 定理スタイルと数式設定 ---
\theoremstyle{definition} % 定義スタイル
\numberwithin{equation}{section} % 式番号をサブセクション単位でリセット

% --- hyperrefの設定 ---
\hypersetup{
  setpagesize = false,
  bookmarks = true,
  bookmarksdepth = tocdepth,
  bookmarksnumbered = true,
  colorlinks = false,
  pdftitle = {}, % PDFタイトル
  pdfsubject = {}, % PDFサブジェクト
  pdfauthor = {}, % PDF作者
  pdfkeywords = {} % PDFキーワード
}

% --- siunitxの設定 ---
\sisetup{
  table-format = 1.5, % 小数点以下の桁数
  table-number-alignment = center, % 数値の中央揃え
}

% --- 透かし画像の設定 --- 
\backgroundsetup{
  scale=0.5,                       % 画像のスケール
  % color=black,                   % 画像の色（透かし用に半透明が推奨）
  opacity=0.2,                   % 透かしの透明度（0が完全透明、1が完全不透明）
  angle=0,                       % 画像の角度
  position = current page.south east,  % ページの右下
  hshift=-6cm, % 右方向へのシフト（負の値で内側に移動）
  vshift=5cm,  % 上方向へのシフト（正の値で内側に移動）
  contents={\includegraphics{./fig/appilogo-circular-full.png}} % 画像のパス
}

% --- その他の設定 ---
\allowdisplaybreaks % 数式の途中改ページ許可
\newcolumntype{t}{!{\vrule width 0.1pt}} % 新しいカラムタイプ
\newcolumntype{b}{!{\vrule width 1.5pt}} % 太いカラム
\UseTblrLibrary{amsmath, booktabs, counter, diagbox, functional, hook, html, nameref, siunitx, varwidth, zref} % tabularrayのライブラリ
\setlength{\columnseprule}{0.4pt} % カラム区切り線の太さ
\captionsetup[figure]{font = bf} % 図のキャプションの太字設定
\captionsetup[table]{font = bf} % 表のキャプションの太字設定
\captionsetup[lstlisting]{font = bf} % コードのキャプションの太字設定
\captionsetup[subfigure]{font = bf, labelformat = simple} % サブ図のキャプション設定
\setcounter{secnumdepth}{4} % セクションの深さ設定
\newcolumntype{d}{D{.}{.}{5}} % 数値のカラム
\newcolumntype{M}[1]{>{\centering\arraybackslash}m{#1}} % センター揃えのカラム
\DeclareMathOperator{\diag}{diag}
\everymath{\displaystyle} % 数式のスタイル
\newcommand{\inner}[2]{\left\langle #1, #2 \right\rangle}
\renewcommand{\figurename}{図}
\renewcommand{\i}{\mathrm{i}} % 複素数単位i
\renewcommand{\laplacian}{\grad^2} % ラプラシアンの記号
\renewcommand{\thesubfigure}{(\alph{subfigure})} % サブ図の番号形式
\newcommand{\m}[3]{\multicolumn{#1}{#2}{#3}} % マルチカラムのショートカット
\renewcommand{\r}[1]{\mathrm{#1}} % mathrmのショートカット
\newcommand{\e}{\mathrm{e}} % 自然対数の底e
\newcommand{\Ef}{E_{\mathrm{F}}} % フェルミエネルギー
\renewcommand{\c}{\si{\degreeCelsius}} % 摂氏記号
\renewcommand{\d}{\r{d}} % d記号
\renewcommand{\t}[1]{\texttt{#1}} % タイプライタフォント
\newcommand{\kb}{k_{\mathrm{B}}} % ボルツマン定数
% \renewcommand{\phi}{\varphi} % ϕをφに変更
\renewcommand{\epsilon}{\varepsilon}
\newcommand{\fullref}[1]{\textbf{\ref{#1} \nameref{#1}}}
\newcommand{\reff}[1]{\textbf{図\ref{#1}}} % 図参照のショートカット
\newcommand{\reft}[1]{\textbf{表\ref{#1}}} % 表参照のショートカット
\newcommand{\refe}[1]{\textbf{式\eqref{#1}}} % 式参照のショートカット
\newcommand{\refp}[1]{\textbf{コード\ref{#1}}} % コード参照のショートカット
\renewcommand{\lstlistingname}{コード} % コードリストの名前
\renewcommand{\theequation}{\thesection.\arabic{equation}} % 式番号の形式
\renewcommand{\footrulewidth}{0.4pt} % フッターの線
\newcommand{\mar}[1]{\textcircled{\scriptsize #1}} % 丸囲み文字
\newcommand{\combination}[2]{{}_{#1} \mathrm{C}_{#2}} % 組み合わせ
\newcommand{\thline}{\noalign{\hrule height 0.1pt}} % 細い横線
\newcommand{\bhline}{\noalign{\hrule height 1.5pt}} % 太い横線

% --- カスタム色定義 ---
\definecolor{burgundy}{rgb}{0.5, 0.0, 0.13} % バーガンディ色
\definecolor{charcoal}{rgb}{0.21, 0.27, 0.31} % チャコール色
\definecolor{forest}{rgb}{0.0, 0.35, 0} % 森の緑色

% --- カスタム定理環境の定義 ---
\newtcbtheorem[number within = chapter]{myexc}{練習問題}{
  fonttitle = \gtfamily\sffamily\bfseries\upshape,
  colframe = forest,
  colback = forest!2!white,
  rightrule = 1pt,
  leftrule = 1pt,
  bottomrule = 2pt,
  colbacktitle = forest,
  theorem style = standard,
  breakable,
  arc = 0pt,
}{exc-ref}
\newtcbtheorem[number within = chapter]{myprop}{命題}{
  fonttitle = \gtfamily\sffamily\bfseries\upshape,
  colframe = blue!50!black,
  colback = blue!50!black!2!white,
  rightrule = 1pt,
  leftrule = 1pt,
  bottomrule = 2pt,
  colbacktitle = blue!50!black,
  theorem style = standard,
  breakable,
  arc = 0pt
}{proposition-ref}
\newtcbtheorem[number within = chapter]{myrem}{注意}{
  fonttitle = \gtfamily\sffamily\bfseries\upshape,
  colframe = yellow!20!black,
  colback = yellow!50,
  rightrule = 1pt,
  leftrule = 1pt,
  bottomrule = 2pt,
  colbacktitle = yellow!20!black,
  theorem style = standard,
  breakable,
  arc = 0pt
}{remark-ref}
\newtcbtheorem[number within = chapter]{myex}{例題}{
  fonttitle = \gtfamily\sffamily\bfseries\upshape,
  colframe = black,
  colback = white,
  rightrule = 1pt,
  leftrule = 1pt,
  bottomrule = 2pt,
  colbacktitle = black,
  theorem style = standard,
  breakable,
  arc = 0pt
}{example-ref}
\newtcbtheorem[number within = chapter]{exc}{Requirement}{myexc}{exc-ref}
\newcommand{\rqref}[1]{{\bfseries\sffamily 練習問題 \ref{exc-ref:#1}}}
\newtcbtheorem[number within = chapter]{definition}{Definition}{mydef}{definition-ref}
\newcommand{\dfref}[1]{{\bfseries\sffamily 定義 \ref{definition-ref:#1}}}
\newtcbtheorem[number within = chapter]{prop}{命題}{myprop}{proposition-ref}
\newcommand{\prref}[1]{{\bfseries\sffamily 命題 \ref{proposition-ref:#1}}}
\newtcbtheorem[number within = chapter]{rem}{注意}{myrem}{remark-ref}
\newcommand{\rmref}[1]{{\bfseries\sffamily 注意 \ref{remark-ref:#1}}}
\newtcbtheorem[number within = chapter]{ex}{例題}{myex}{example-ref}
\newcommand{\exref}[1]{{\bfseries\sffamily 例題 \ref{example-ref:#1}}}
% --- 再定義コマンド ---
% \mathtoolsset{showonlyrefs=true} % 必要な式番号のみ表示
\pagestyle{fancy} % ヘッダー・フッターのスタイル設定
\chead{応用量子物性講義ノート} % 中央ヘッダー
% \rhead{}
\fancyhead[R]{\rightmark}
\renewcommand{\sectionmark}[1]{\markright{\thesection\ #1}}
\cfoot{\thepage} % 中央フッターにページ番号
\lhead{}
\rfoot{Yuto Masuda and Haruki Aoki} % 右フッターに名前
\setcounter{tocdepth}{4} % 目次の深さ
\makeatletter
\@addtoreset{equation}{section} % サブセクションごとに式番号をリセット
\makeatother

% --- メタ情報 ---
\title{応用量子物性講義ノート}
\date{更新日\today}
\author{Yuto Masuda and Haruki Aoki}

\begin{document}
  量子力学における散乱では，散乱の前後で位相が変化する．これを\textbf{位相シフト}(Phase shift)とよぶ．
  以下の1次元の例で位相シフトを確認する．
  ポテンシャル$V(x)$の壁が，
  \begin{align}
    V(x) = 
    \begin{dcases}
      0\ &(x < 0) \\
      V_0\ &(x \geq 0)
    \end{dcases}
  \end{align}
  のように書かれて，その壁に左から入射する粒子を考える．
  入射粒子のエネルギーは$E<V_0$とする．Schrödinger方程式を解くことにより波動関数は，
  \begin{align}
    \psi(x) = 
    \begin{dcases}
      \e^{\i kx} + \frac{k -\i\alpha}{k + \i\alpha} \e^{-\i kx} & x < 0 \label{1Dscattering} \\
      \frac{2k}{k + \i\alpha}\e^{-\alpha x} & x \geq 0
    \end{dcases}
  \end{align}
  となる．ただし，$k \coloneqq \frac{\sqrt{2mE}}{\hbar}$，$\alpha \coloneqq \frac{\sqrt{2m(V - E)}}{\hbar}$とおいた．
  \refe{1Dscattering}の第1式において，第1項は入射波を，第2項は反射波を表している．
  反射波は，$\delta_k \coloneqq \arg\qty(k + \i\alpha)$とすると，
  波動関数は，
  \begin{align}
    \psi(x) &= \e^{\i kx} + \e^{-2 \i \delta_k}\e^{-\i kx} \\ 
    &= \e^{2\i\delta_k}\qty(\e^{\i jx}\e^{2\i\delta_k} + \e^{-\i kx})
  \end{align}
  となる．この$\delta_k$は位相シフトと呼ばれ，量子力学的散乱を特徴づけるパラメータである．
  \par
  位相シフトの現象を3次元に拡張する．
  今回は全反射が起こるとする．
  \refe{PartialWave}において$r\to\infty$とすると，
  \begin{align}
    j_l(kr) \to \frac{\exp\qty{\i\qty(kr - \frac{l\pi}{2})} - \exp\qty{-\i\qty(kr - \frac{l\pi}{2})}}{2\i kr}
  \end{align}
  であるため，
  \begin{align}
    \psi(\bm{r}) = \sum_{l = 0}^{\infty} \frac{2l + 1}{2\i kr}\qty[\i^l \qty(\exp\qty{\i\qty(kr - \frac{l\pi}{2})} - \exp\qty{-\i\qty(kr - \frac{l\pi}{2})}) + (2\i k) a_l \e^{\i kr}]P_l(\cos\theta)
  \end{align}
  を得る．さらに，
  \begin{align}
    \i^l \qty[\exp\qty{\i\qty(kr - \frac{l\pi}{2})} - \exp\qty{-\i\qty(kr - \frac{l\pi}{2})}] &= \i^l \exp\qty(-\i\frac{l\pi}{2})\qty(\e^{\i kr} - \e^{-\i kr} e^{\i l\pi}) \\
    &= 1\cdot\qty(\e^{\i kr} - (-1)^l) \e^{-\i kr}
  \end{align}
  なる関係を用いると波動関数は，
  \begin{align}
    \psi(\bm{r}) = \sum_{l = 0}^{\infty} \qty[\qty(1 + 2\i ka_l)\e^{\i kr} - (-1)^l \e^{-\i kr}]P_l(\cos\theta)\label{SCPartialWave}
  \end{align}
  と書ける．
  第1項は外向き球面波，第2項は内向き球面波を表す．この散乱は全反射であるため入射波と反射波の振幅は等しい．
  つまり，
  \begin{align}
    \abs{1 + 2\i ka_l} = 1
  \end{align}
  が成り立つ．よって，散乱による位相のずれを$\delta_l$とおくと$a_l$は，
  \begin{align}
    &1 + 2\i ka_l = \e^{2\i \delta_l} \\ 
    &\Leftrightarrow a_l = \frac{1}{2\i k}(\e^{2\i \delta_l} - 1) \\
    &\Leftrightarrow a_l = \frac{1}{2\i k}(\cos^2\delta_l + 2\i\sin\delta_l\cos\delta_l - \sin^2\delta_l - 1) \\
    &\Leftrightarrow a_l = \frac{1}{2\i k}(1 - \sin^2\delta_l + 2\i\sin\delta_l\cos\delta_l - \sin^2\delta_l - 1) \\
    &\Leftrightarrow a_l = \frac{1}{2\i k}(2\i\sin\delta_l\cos\delta_l - 2\sin^2\delta_l) \\
    &\Leftrightarrow a_l = \frac{1}{k}(\cos\delta_l + \i\sin\delta_l)\sin\delta_l\\
    &\Leftrightarrow a_l = \frac{1}{k}\e^{\i \delta_l} \sin \delta_l
  \end{align}
  である．
  全断面積$\sigma^{\r{tot}}$は，
  \begin{align}
    \sigma^{\r{tot}} &= 4\pi \sum_{l=0}^{\infty} (2l + 1)\abs{a_l}^2 \\
    &= \frac{4\pi}{k^2} \sum_{l=0}^{\infty} (2l + 1)\sin^2\delta_l 
  \end{align}
  である．
  散乱振幅$f(\theta)$は，
  \begin{align}
    f(\theta) &= \sum_{l = 0}^{\infty} (2l + 1)a_lP_l(\cos\theta) \\
    &= \frac{1}{k} \sum_{l = 0}^{\infty}(2l + 1)\e^{\i\delta_l}\sin\delta_lP_l(\cos\theta)
  \end{align}
  である．
  特に$\theta = 0$のとき，
  \begin{align}
    f(0) &= \frac{1}{k} \sum_{l = 0}^{\infty} (2l + 1) (\cos\delta_l + \i \sin\delta_l)\sin\delta_l P_l(1) \\ 
    &= \frac{1}{k} \sum_{l = 0}^{\infty} (2l + 1) (\cos\delta_l + \i \sin\delta_l)\sin\delta_l
  \end{align}
  が成り立つ．したがって以下の\textbf{光学定理}が成り立つ．
  \begin{itembox}[l]{光学定理}
    \begin{align}
      \sigma^{\r{tot}} = \frac{4\pi}{k} \Im f(0)
    \end{align}
  \end{itembox}
  これは全断面積が前方散乱の散乱振幅からわかることを示している\footnote{
    砂川，散乱の量子論，「光学定理は，前方散乱によって，入射波の強度が減少した分だけ，四方に散乱されるという，まことに当然なことを述べているのである．」 
  }．
  \begin{myex}{半径$a$の剛体球による散乱(再考)}{}
    散乱体のポテンシャルを，
    \begin{align}
      V(r) =
      \begin{dcases}
        \infty & r \leq a  \\
        0 & r \geq a
      \end{dcases}
    \end{align}
    とする．低エネルギー散乱($ka \ll 1$)とする．このとき，散乱の影響を受けるのはほぼ$l = 0$のみである．
    \refe{SCPartialWave}において$l = 0$とすることで，
    \begin{align}
      \psi(\bm{r}) = \frac{1}{2\i kr}\qty[(1 + 2\i k a_0)\e^{\i kr} - \e^{-\i kr}] P_0(\cos\theta)
    \end{align}
    を得る．位相シフトを考慮し$\e^{2\i \delta_0} = 1 + 2\i ka_0$とおく．境界条件より
    \begin{align}
      \psi(a) = 0
    \end{align}
    であるから
    \begin{align}
      \e^{2\i \delta_0}\e^{\i ka} - \e^{-\i ka} = 0
    \end{align}
    が成り立つ．よって位相シフトは，
    \begin{align}
      \delta_0 = -ka
    \end{align}
    である．全断面積は，
    \begin{align}
      \sigma^{\r{tot}} &= \frac{4\pi}{k^2}\sin^2\delta_0 \\
      &\simeq \frac{4\pi}{k^2}\delta_0 \\
      &= 4\pi a^2
    \end{align}
    である．これは古典力学における剛体級の散乱$\sigma^{\r{tot}}=\pi a^2$の4倍の値である．波長$\lambda=\frac{2\pi}{k}$が散乱体の半径$a$より十分に大きいため，
    回折によって剛体球を取り囲み，球の表面積を疑似的に増加させたためと説明できる．
  \end{myex}
  \begin{myex}{位相シフトの問題}{}
    散乱体から十分遠方での境界条件を満たし，ポテンシャル外における部分波展開した波動関数として，
    \begin{align}
      \psi(\bm{r}) = \sum_{l=0}^{\infty} \i^l \frac{2l+1}{2}\qty[\e^{2\i\delta_l} h_l^{(1)} (kr) + h_l^{(2)} (kr)] P_l(\cos\theta)
    \end{align}
    を用いて，半径$a$の剛体球，
    \begin{align}
      V(r) = 
      \begin{dcases}
        \infty & r\leq a \\
        0 & r > a 
      \end{dcases}
    \end{align}
    による散乱を考える．特に，$ka$が小さい低エネルギー散乱において，$l$の増大に伴ってその寄与が散乱に対してどの程度小さくなるか
    調べる．ここで，
    \begin{align}
      h_l^{(1)}(kr) = j_l(kr) + \i y_l(kr),\ h_l^{(2)}(kr) = j_l(kr) - \i y_l(kr)
    \end{align}
    は球Hankel関数であり，$j_l(kr)$は球Bessel関数，$y_l(kr)$は球Neumann関数である．
    \begin{enumerate}
      \item 波動関数の境界条件を考えることで，散乱による位相シフト$\tan\delta_l$を$j_l(ka)$と$y_l(ka)$を用いて表せ．
      \item 低エネルギー散乱($ka \ll 1$)を考える．このとき，
        \begin{align}
          j_l(ka) \to \frac{(ka)^l}{(2l+1)!!},\ y_l(ka) \to -\frac{(2l-1)!!}{(ka)^{l+1}}\label{LimitOfBessel}
        \end{align}
        \begin{align}
          (2l+1)!! = (2l+1)(2l-1)\cdots5\cdot3\cdot1 \\
          (2l-1)!! = (2l-1)(2l-3)\cdots3\cdot1\cdot1
        \end{align}
        であることを用いて，$l = 0, 1, 2$についての位相シフト$\tan\delta_l$を求めよ．得られる$\tan\delta_0$，$\tan\delta_1$，$\tan\delta_2$は
        $l$の増加と共に位相シフトが急激に小さくなり，$ka$が小さい時，$l = 0$だけ考えれば十分であることを示すものである．
    \end{enumerate}
    \tcblower
    \begin{enumerate}
      \item 波動関数は剛体球の中に侵入できないため$r=a$で$\psi=0$である．よって，
        \begin{align}
          \e^{2\i \delta_l} h_l^{(1)} (ka) + h_l^{(2)}(ka) = 0 \\
          \e^{2\i \delta_l}\qty(j_l(ka) + \i y_l(ka)) + \qty(j_l(ka) - \i y_l(ka)) = 0 \\
          \qty(\e^{2\i \delta_l} + 1)j_l(ka) + \i\qty(\e^{2\i \delta_l} - 1)y_l(ka) = 0 \\
          2\cos\delta_l j_l(ka) -2 \sin\delta_l y_l(ka) = 0
        \end{align}
        となり，位相シフトとして，
        \begin{align}
          \tan\delta_l = \frac{j_l(ka)}{y_l(ka)}\label{phaseshift}
        \end{align}
        が得られる．
      \item \refe{LimitOfBessel}を\refe{phaseshift}に代入すると，
        \begin{align}
          \tan\delta_l &\to \frac{(ka)^l}{(2l+1)!!} \qty(- \frac{(ka)^{l + 1}}{(2l -1)!!}) \\
          &= -\frac{(ka)^{2l + 1}}{(2l + 1)!!(2l - 1)!!} \\
        \end{align}
        となる．
        $l = 0, 1, 2$を代入する．ここで，$(n - 2)!! = \frac{n!}{n}$であることに注意すると，
        \begin{align}
          \tan\delta_0 &= -ka \\
          \tan\delta_1 &= -\frac{1}{3}(ka)^3 \\
          \tan\delta_2 &= -\frac{1}{45}(ka)^5
        \end{align}
        となる．
    \end{enumerate}
  \end{myex}
\end{document}